% 旧中心商特異性論文VI 第15.3・16節の現行TeX回復
% 有限中心・Azumaya・中心chart gluingは現行証明を参照し、PBW/CY等の基礎入力を明示する。
\section{一次元表現の非被約性と正確なtails正則判別}
\label{sec:quotient-exact-tails}
この節は旧論文VIの第15.3節・第16節の証明を現在有効な順序で整理する。
結論は$h=1$の指定した一混合項族に限る。
全21混合項の一般判別式や任意の量子パラメータ$h$の判別式を
主張しない。有限中心・Azumaya locusは第\ref{sec:quotient-center-azumaya}節、
中心chartの貼合せは第\ref{sec:quotient-chart-gluing}節の証明を用いる。
PBW、周囲のAS正則性、$R_1$のCalabi--Yau性は
第\ref{sec:quotient-pbw-cy}節で証明した。
斉次Sklyaninの楕円境界の記述は既知入力として区別する。

$R_1=\CC\langle X,Y,Z\rangle/
([X,Y]-Z^2,[Y,Z]-(X^2+1),[Z,X]-Y^2)$、
$G_1=\Omega_1(X,Y,Z)+X$とし、
$D_\beta=R_1/(G_1+\beta)$と置く。

\begin{proposition}[一次元表現を検出する中心値]
\label{prop:tails-one-dimensional}
$D_\beta$に一次元表現があるための必要十分条件は
$9\beta^2+4=0$である。
$s^2=-1$、$\beta=-2s/3$なら、
唯一の一次元表現の局所スキームは
\[
 \operatorname{Rep}_1(D_{-2s/3})
 \simeq\operatorname{Spec}\frac{\CC[y,z]}{(y^2,yz,z^2)}
\]
であり、長さ3、接空間次元2である。
全$\beta$にわたる一次元表現のスキーム論的像の
定義イデアルは$((9\beta^2+4)^2)$である。
\end{proposition}
\begin{proof}
可換化を取ると
$T=\CC[x,y,z]/(x^2+1,y^2,z^2)$となり、
$G_1$の像は$\frac23x(1-yz)$である。
$x^2+1$の二根$s=\pm i$の各因子で
$G_1+\beta=\beta+\frac23s-\frac23syz$となる。
前半はこの式の単元性を調べれば得られる。
さらに$\beta=-2s/3$では$yz=0$が追加され、
所望の局所環になる。
全因子にわたる中心写像は
$\beta=-\frac23x(1-yz)$である。
$x^2=-1$および$(yz)^2=0$から
$9\beta^2+4=8yz$であり、
各因子で$yz\ne0$なので像の核は
$((9\beta^2+4)^2)$となる。
\end{proof}

\begin{lemma}[有限中心を持つ中心正則商のCY双対性]
\label{lem:tails-central-cy}
$R$を$\CC$上の$d$次元Calabi--Yau代数、
$w\in Z(R)$を左右正則元とし$D=R/(w)$とする。
$D$のhomological smoothnessを仮定しなくても、
双対化複体には
\[
 \mathbf R\operatorname{Hom}_{D^e}(D,D^e)
 \simeq D[-(d-1)]
\]
という同型がある。
さらに$D$が有限生成可換中心部分代数上有限で
有限大域次元を持てば、
$D$はhomologically smoothで$(d-1)$-Calabi--Yauである。
\end{lemma}
\begin{proof}
$R^e=R\otimes R^{\mathrm{op}}$内の中心正則列
$w\otimes1,1\otimes w$のKoszul分解から
$\mathbf R\operatorname{Hom}_{R^e}(D^e,R^e)
 \simeq D^e[-2]$を得る。
一方、$0\to R\xrightarrow{w}R\to D\to0$と
$\mathbf R\operatorname{Hom}_{R^e}(R,R^e)\simeq R[-d]$
から
$\mathbf R\operatorname{Hom}_{R^e}(D,R^e)\simeq D[-d-1]$。
導来随伴で二つを比較すれば最初の表示が従う。
有限型可換中心部分代数$C$上有限なら、$D$と$D^{\mathrm{op}}$は
affine PIかつ左右noetherianで、$D^e$も$C\otimes C$上有限だからnoetherian。
有限大域次元を仮定した場合、
\cite[Lemmas~3.4, 3.5]{RSVBertini}から
$D^{\mathrm{op}}$と$D^e$も有限大域次元を持つ。
ここで基礎体$\CC$の代数閉性とaffine PI性を使った。
従って有限生成$D^e$加群$D$はperfectでsmoothとなる。
非可換中心商が無条件にsmoothになるとは主張していない。
\end{proof}

\begin{theorem}[二つの真に悪いaffine中心値]
\label{thm:tails-bad-affine}
命題\ref{prop:quotient-constant-cy}と
第\ref{sec:quotient-center-azumaya}節により、
$R_1$は三次元Calabi--Yauで中心上有限、$G_1+\beta$は中心正則元である。
このとき
\[
 \beta=\pm\frac{2i}{3}
 \quad\Longrightarrow\quad \operatorname{gldim}D_\beta=\infty.
\]
\end{theorem}
\begin{proof}
もし大域次元が有限なら、
補題\ref{lem:tails-central-cy}から$D_\beta$は
二次元Calabi--Yauとなる。

Bocklandt--Galluzzi--\allowbreak Vaccarinoの単純表現の局所構造定理
\cite[Theorems 6.3, 6.6 and Corollary 6.8]{BGV}
によれば、その一次元単純表現の点で
$\operatorname{Rep}_1(D_\beta)$は滑らかである必要がある。
しかし命題\ref{prop:tails-one-dimensional}の局所環は
零次元・接空間次元2であり非正則だから矛盾する。
\end{proof}

\begin{lemma}[有限射影行列因子化によるnonperfectness検出]
\label{lem:tails-finite-mf}
$R$を左右noetherian、$\operatorname{gldim}R\le d<\infty$、
$w\in Z(R)$を正則元とし、$D=R/(w)$と置く。
有限生成$D$加群の$(d-1)$次syzygyは
有限生成射影$R$加群の行列因子化のcokernelとして得られる。
さらにそのcokernelが有限射影次元なら
当該行列因子化はhomotopy縮約可能である。
従って$\operatorname{gldim}D=\infty$なら
非縮約な有限射影行列因子化が存在する。
\end{lemma}
\begin{proof}
$\operatorname{pd}_R D=1$なので、有限自由$D$分解を
$(d-1)$回ずらしたsyzygy $N$は
$R$上の長さ1の有限射影分解
$0\to P_1\xrightarrow{\phi}P_0\to N\to0$を持つ。
$wN=0$から$\phi\psi=w$を満たす$\psi$を得て
単射性から$\psi\phi=w$となる。
$D$上で有限射影次元なら補題\ref{lem:gluing-contraction}で縮約できる。
全ての有限生成加群の射影次元が有限なら同じ構成で
上界$d-1$を一様に与え、最後の主張が従う。
\end{proof}

\begin{lemma}[中心値の二重分岐はnonperfectnessを保存]
\label{lem:tails-double-zero}
$R$は左右noetherianで有限大域次元を持ち、
$G+\beta_0\in Z(R)$は正則で
$\operatorname{gldim}R/(G+\beta_0)=\infty$とする。
$s_0\in\CC^\times$、$b(s)\in\CC[s,s^{-1}]$が
\[
 b(s)-\beta_0=a(s)c(s),\quad
 a(s_0)=c(s_0)=0
\]
を満たし、$G+b(s)$は$R[s,s^{-1}]$上正則とする。
このとき
$T=R[s,s^{-1}]/(G+b(s))$は有限生成非perfect加群を持つ。
\end{lemma}
\begin{proof}
補題\ref{lem:tails-finite-mf}より
$G+\beta_0$の非縮約な有限射影行列因子化
$(P,d_0)$を取る。
これを$R[s,s^{-1}]$に拡大し、
中心的階数1の因子化$(a,c)$とsuper tensorを取る。
全微分$d$は$d^2=G+\beta_0+ac=G+b(s)$を満たす。
そのcokernelがperfectなら前補題の縮約の議論により
$d$は縮約できる。
$s=s_0$へ特殊化すると$(a,c)$は零微分となり、
$d_0\oplus d_0[1]$が縮約する。
直和因子に射影すれば$d_0$も縮約して矛盾。
\end{proof}

\begin{lemma}[tailsから中心chart、有限\'etale基底変更へ]
\label{lem:tails-chart-perfect}
$A$をnoetherian connected graded代数とし、
$u\in Z(A)_1$を正則元とする。
$\operatorname{gldim}\operatorname{qgr}A<\infty$なら
$Q_0=A[u^{-1}]_0$の任意の有限生成加群はperfectである。
この性質は中心局所化、
および有限\'etale可換中心拡大$Q\to T$にも伝わる。
\end{lemma}
\begin{proof}
$Q_0$の有限生成加群$M$の有限個の生成元から、
$M[u,u^{-1}]$内に有限生成graded $A$-latticeを取れる。
$\operatorname{qgr}A$の有限大域次元とtwistsの生成により、
その有界導来圏は$\pi A(m)$で厚く生成される。
中心局所化後、twistsは$Q_0$上自由だから$M$はperfect。
中心局所化には有限生成latticeが存在しperfect性を保つ。
有限\'etale拡大では$M$は
$M\otimes_QT$の$T$-直和因子であり、
separability idempotentによりperfect性が降下する。
\end{proof}

\begin{proposition}[曲面の判別式：中心chartからの導出]
\label{prop:tails-surface-discriminant}
$C_\alpha=A_\alpha/(v-u)$とし、
$B_\alpha,A_\alpha$の定義を
第\ref{sec:quotient-exact-tails}節の後半の定理に置く。
第\ref{sec:quotient-pbw-cy}節のPBW・AS正則性・CY性と、
境界$\operatorname{qgr}(C_\alpha/(u))\simeq\operatorname{coh}E_1$
の既知の記述（$E_1$は滑らかな楕円曲線）を用いる。
貼合せ自体を追加仮定とはしない。
すると
\[
 \operatorname{gldim}\operatorname{qgr}C_\alpha=
 \begin{cases}2,&(\alpha^3+9)(28\alpha^3+243)\ne0,\\
 \infty,&(\alpha^3+9)(28\alpha^3+243)=0.\end{cases}
\]
\end{proposition}
\begin{proof}
$\alpha\ne0$では中心chartは
$R_1/(G_1+\beta)$、$\beta^2\alpha^3=4$に同型。
中心曲面判別式とAzumaya性から
$\gamma^2\notin\{0,-1,-28/27\}$では
$\operatorname{gldim}R_1/(g-\gamma)=2$である
（$g=3G_1/2$、$\gamma=-3\beta/2$）。
$\gamma^2=-28/27$では前節の中心Azumaya点における
特異曲面から無限大域次元、
$\beta=\pm2i/3$では
定理\ref{thm:tails-bad-affine}から無限大域次元。
後二者は補題\ref{lem:tails-chart-perfect}の逆向き対偶により
tailsにも障害となる。
よい値では楕円境界は有限大域次元を持ち、
そのcoherent対象は構造対象のtwistsで厚く生成される。
実際、twistsの和による解消を大域次元より一つ長く切り、
最後のYoneda拡大の消滅から各対象を有限複体の直和因子として得る。
命題\ref{prop:gluing-one-chart}で全coherent対象のperfect性を得て、
AS--Gorenstein次元3、index 1の双対化複体に
補題\ref{lem:gluing-serre-dimension}を適用するとtails大域次元は正確に2となる。
$\alpha=0$では$u$-chartが$S_1/(\Omega_1+2)$となる。
Euler導分$\delta(\Omega_1)=3\Omega_1$に
補題\ref{lem:gluing-moving-central-value}を適用すると、
$P(-2)=-6\ne0$からchartの大域次元は高々2となる。
同じ命題とSerre双対性が適用できる。
\end{proof}

\begin{theorem}[一混合項族の三次三様体と曲面の正確な判別式]
\label{thm:quotient-exact-tails}
\[
 B_\alpha=
 \frac{\CC[u,v]\langle x,y,z\rangle}
 {([x,y]-z^2,\ [y,z]-(x^2+\alpha uv),\
   [z,x]-y^2)},\qquad
 A_\alpha=B_\alpha/(W_\alpha),
\]
\[
 W_\alpha=\frac{x^3+y^3+z^3}{3}
 -\frac{xyz+xzy+yxz+yzx+zxy+zyx}{9}
 +\alpha uvx+u^3+v^3,\quad C_\alpha=A_\alpha/(v-u)
\]
とする。系\ref{cor:quotient-ambient-pbw}と
命題\ref{prop:quotient-constant-cy}により、
$B_\alpha$はKoszul AS正則（次元5、index 5）、
$W_\alpha$は中心正則、$A_\alpha$は$\CC[u,v]$上自由で、
$R_1$は三次元Calabi--Yauである。
斉次Sklyaninの既知の記述
$\operatorname{qgr}(S_1/(\Omega_1))\simeq\operatorname{coh}E_1$
（$E_1$は滑らかな楕円曲線）を用いる。
最後の境界の代数的表示は
\cite[Lemma~2.19]{WWYSklyanin}のtwisted homogeneous coordinate ringである。
中心chartの貼合せは第\ref{sec:quotient-chart-gluing}節で証明済みである。
$\Delta(\alpha)=(\alpha^3+9)(28\alpha^3+243)$と置くと
\[
 \bigl(\operatorname{gldim}\operatorname{qgr}C_\alpha,\,
 \operatorname{gldim}\operatorname{qgr}A_\alpha\bigr)
 =
 \begin{cases}(2,3),&\Delta(\alpha)\ne0,\\
 (\infty,\infty),&\Delta(\alpha)=0.
 \end{cases}
\]
従って全ての$\alpha\in\mathbb Q$は正則であり、
$\alpha=1$は有理係数の具体的な本質混合例を与える。
\end{theorem}
\begin{proof}
曲面は前命題による。
三次三様体のよい値では、
中心導分$u\partial_u-v\partial_v$の
$\delta(W_\alpha)=3(u^3-v^3)$と、
中心切断$v=\omega u$（$\omega^3=1$）を用いる。
この導分は$uv$を固定するため、関係式を保つ真の導分である。
$\delta(W_\alpha)=-3\prod_{\omega^3=1}(v-\omega u)$であり、
各因子は$\CC[u,v]$上の自由性から$A_\alpha$上正則。
自己同型$(u,v)\mapsto(\zeta u,\zeta^{-1}v)$（$\zeta^3=1$）は
周囲の関係式と$W_\alpha$を保ち、三つの切断を巡回的に交換する。
各切断は前命題の$C_\alpha$と同型なので、
命題\ref{prop:gluing-derivation}から全coherent対象のperfect性を得る。
$A_\alpha$のAS--Gorenstein次元4、index 2を
補題\ref{lem:gluing-serre-dimension}に適用すれば、
$\operatorname{gldim}\operatorname{qgr}A_\alpha=3$となる。

悪い値では$\alpha\ne0$である。
$u,v$を逆にした次数零chartで$t=v/u$とし、
中心に$s^2=\alpha t$を添加する。
$s,2s$が可逆なのでこれは有限\'etale拡大となり、
$x=sX,y=sY,z=sZ$とすると
\[
 T\simeq R_1[s,s^{-1}]/(G_1+b_\alpha(s)),
 \qquad b_\alpha(s)=s^{-3}+\frac{s^3}{\alpha^3}.
\]
$s_0^2=\alpha$とし$\beta_0=2/s_0^3$と置く。
$\beta_0^2\alpha^3=4$であり、
$\Delta(\alpha)=0$なら
$\beta_0^2=-4/9$または$-112/243$。
前者は定理\ref{thm:tails-bad-affine}、
後者は前節のAzumaya局所判定から
$\operatorname{gldim}R_1/(G_1+\beta_0)=\infty$。
一方
\[
 b_\alpha(s)-\beta_0
 =\frac{(s^3-s_0^3)^2}{s_0^6s^3}.
\]
ここで$a(s)=s^3-s_0^3$、
$c(s)=a(s)/(s_0^6s^3)$はともに$s=s_0$で零。
補題\ref{lem:tails-double-zero}により$T$に
有限生成非perfect加群が存在する。
もし$\operatorname{gldim}\operatorname{qgr}A_\alpha<\infty$なら
補題\ref{lem:tails-chart-perfect}により$T$の全有限生成加群はperfectで、
矛盾する。これで必要性を得る。
最後に$\alpha\in\mathbb Q$なら
$\alpha^3=-9,-243/28$は三進付値により不可能。
\end{proof}

\begin{remark}[量子的判別式と古典・点加群の違い]
$\alpha^3=-243/28$では古典三次三様体と
量子的楕円境界は滑らかなまま、tailsのみ非正則となる。
さらに非零$\alpha$で$\alpha^3\ne-9$の範囲では
正則例と新しい非正則例の定常点加群スキームは
六点を台とする長さ4の局所環
$\CC[Y,Z]/(Y^2,Z^2)$を持つ。
これは埋込を含む点スキーム全体の同型を主張しない。
\end{remark}


\begin{corollary}[全affine中心値の必要十分判別式]
\label{cor:tails-all-affine}
全ての$\beta\in\CC$について
\[
 \operatorname{gldim}R_1/(G_1+\beta)
 =\begin{cases}
 2,&(9\beta^2+4)(243\beta^2+112)\ne0,\\
 \infty,&(9\beta^2+4)(243\beta^2+112)=0.
 \end{cases}
\]
\end{corollary}
\begin{proof}
$\beta=0$は系\ref{cor:zero}による。
それ以外は$g=-3\beta/2$、$H=9\beta^2/4$として
系\ref{cor:quotient-azumaya-fibers}の滑らかな中心ファイバーと特異Azumaya点の判定を使う。
$9\beta^2+4=0$の場合は定理\ref{thm:tails-bad-affine}による。
\end{proof}
